\section{非微扰结果}
\label{sec:nonper}
\begin{figure}[htbp]
    \centering
    \includegraphics[width=0.6\textwidth]{images/zvl.pdf}
    \caption{\itshape $Z_{V^L}$ 作为 $\mu^2 a^2$ 的函数。由于格点在空间与时间方向不对称，用局域流的
时间分量（$V_0$）与三个方向平均的空间分量（$V_k$）所得的结果分别给出。虚线是来自 BSPT 的
$Z^{{\rm BSPT}}_{V^L}$，曲线是 $16^3 \times 32$ 体积上 BDPT 的结果，见第 \ref{sec:per} 节，
见第 \ref{sec:per} 节。直线（连续线）是用 Ward 恒等式方法 \cite{wi,gribov} 得到的结果。}
    \label{fig:zvl}
\end{figure}
\begin{figure}[htbp]
    \centering
    \includegraphics[width=0.6\textwidth]{images/zs.pdf}
    \caption{\itshape $Z_S$ 作为 $\mu^2 a^2$ 的函数。我们给出用夸克波函数重正化的两种可能定义
（记作 $Z_\psi$，式 (\ref{eq:zpsil})；及 $Z^\prime_\psi$，式 (\ref{eq:zpsi1})）所得的重正化常数。虚线是
无限体积极限下 boosted 标准微扰论的 $Z^{{\rm BSPT}}_S$，实线是 BDPT 的结果。}
    \label{fig:zs}
\end{figure}
\begin{figure}[htbp]
    \centering
    \includegraphics[width=0.6\textwidth]{images/za.pdf}
    \caption{\itshape $Z_{A^L}$ 作为 $\mu^2 a^2$ 的函数。我们分别给出用时间分量 $A_0$ 与空间分量
$A_k$ 得到的重正化常数。虚线是 boosted 标准微扰论的 $Z^{{\rm BSPT}}_A$，实线来自 BDPT。直线（连续线）
是用 Ward 恒等式方法 \cite{wi,gribov} 得到的结果。}
    \label{fig:za}
\end{figure}
\begin{figure}[htbp]
    \centering
    \includegraphics[width=0.6\textwidth]{images/zp.pdf}
    \caption{\itshape $Z_P$ 作为 $\mu^2 a^2$ 的函数。我们给出用夸克波函数重正化的两种可能定义
（记作 $Z_\psi$，式 (\ref{eq:zpsil})；及 $Z^\prime_\psi$，式 (\ref{eq:zpsi1})）所得的重正化常数。虚线是
无限体积极限下 boosted 标准微扰论的 $Z^{{\rm BSPT}}_P$，实线是 BDPT 的结果。
}
    \label{fig:zp}
\end{figure}
\begin{figure}[htbp]
    \centering
    \includegraphics[width=0.6\textwidth]{images/zsszp.pdf}
    \caption{\itshape{$Z_S/Z_P$ 作为 $\mu^2 a^2$ 的函数。虚线对应 BSPT，实线来自用 Ward 恒等式方法
\cite{wi} 对该比值的测定。
}}
    \label{fig:zsszp}
\end{figure}

本节给出我们对本文所提非微扰重正化流程数值研究的主要结果，以及这些结果与微扰论的比较。

我们在 $\beta=6.0$ 的 $16^3 \times 32$ 格点上生成 $36$ 个独立规范场组态进行模拟。夸克传播子只对单一夸克质量值
（$am_q \sim 0.07$）计算，对应于跃迁参数 $K=0.1425$。所有格林函数都在格点朗道规范中计算，该规范通过极小化泛函
\beq
{\rm tr}\Bigl[ \sum_{\mu=1}^{4}\Bigl(U_{\mu}(x)+
U^{\dagger}_{\mu}(x)\Bigr)\Bigr].
\eeq
得到。来自 Gribov 副本或伪解的可能效应 \cite{giusti} 未予考虑。如下所示，对于那些也能以规范不变方式确定的量，
在固定规范的夸克态上得到的非微扰结果与用 Ward 恒等式方法 \cite{wi} 所得的结果符合得很好。

\noindent {\bf 矢量流：}

图 \ref{fig:zvl} 给出用式 (\ref{eq:zpsi1}) 与 (\ref{eq:gvl}) 得到的局域矢量流重正化常数 $Z_{V^L}$。如预期那样，
在统计误差内 $Z_{V^L}$ 不依赖于标度，直至格点离散化畸变变得重要的大 $\mu^2$ 值处。这是如下等价性的后果：此处用于确定
$Z_{V^L}$ 的方法与局域矢量流的 Ward 恒等式在 $O(\alpha_s a)$ 精度内可以建立等价关系，见第 \ref{sec:casivari} 节。
利用 $\mu^2 a^2 \sim 1$（相应于 $\mu \sim 2$ GeV）处的点，得到 $Z^{{\rm NP}}_{V^L}=0.84(1)$；
与之相比，Ward 恒等式方法给出 $Z^{{\rm WI}}_{V^L}=0.824(2)$，BSPT 给出 $Z^{{\rm BPT}}_{V^L}=0.83$
（在 $\mu^2 a^2 \sim 1$ 处的 $16^3 \times 32$ 格点上为 $0.86-0.87$）。三种方法符合良好。

\noindent {\bf 标量密度：}

出于若干理由，$Z_S$ 是检验我们方法有效性的理想量：它无法用 Ward 恒等式方法确定；它是对数发散且规范依赖的；
并且（与赝标密度和轴矢流不同）它不受 Goldstone 玻色子极点的影响，见第 \ref{sec:idea} 节。

图 \ref{fig:zs} 给出作为重正化标度函数的 $Z_S$。我们分别报告用 $Z_\psi$（式 (\ref{eq:zpsil})）与
$Z^\prime_\psi$（式 (\ref{eq:zpsi1})）所得的结果。注意，在 $\mu^2 a^2$ 不太大、离散化误差较小处，对应于不同
$Z_\psi$ 定义的点彼此符合得很好。比较 BSPT（虚线）与 BDPT（实线）的结果也可看出这一点。我们不预期与微扰论在低
$\mu^2$（高阶效应非常大）或高 $\mu^2$（格点畸变重要，如实线所示）处相符。数值结果遵循理论预期：在
$0.3 \le \mu^2 a^2 \le 1$（$1.1$ GeV $\le \mu \le 2$ GeV）范围内，$Z_S$ 的非微扰测定与 boosted 微扰论预言符合良好。
注意此处的 $Z_S$（$Z_P$）由离壳格林函数定义，依赖于标量（赝标）密度的反常量纲，因而依赖 log($\mu^2 a^2$)。此外，
即使在连续统中，该结果也是规范依赖的——这一点可由一圈微扰论的显式计算证明。

\noindent {\bf 轴矢流：}

局域轴矢流具有矢量流的某些特征。其重正化常数在微扰论所有阶都有限 \cite{curci}，且可由 Ward 恒等式确定。然而轴矢流与
$QCD$ 的准 Goldstone 玻色子耦合，后者在低 $\mu^2$ 处可能有重要贡献。与 $Z_{V^L}$ 类似、应不依赖于 $\mu$ 的 $Z_{A^L}$
作为 $\mu^2$ 的函数示于图 \ref{fig:za}。

 我们把低 $\mu^2$ 处 $Z_{A^L}$ 对 $\mu$ 的强依赖解释为赝标态的非微扰效应。遗憾的是，在非微扰区与大 $\mu$ 区之间看不到
清晰的平台存在的迹象，后者是格点伪迹变得重要的区域\footnote{在更大的 $\beta$ 值处平台最终可能更清晰地显现。}。若没有
来自 Ward 恒等式方法 \cite{wi} 的信息，将难以自信地确定 $Z_{A^L}$。尽管如此，我们的确观察到：在 $\mu^2 a^2 \sim 1$
处（恰好在 DPT 中格点伪迹变大之前），结果接近于用 Ward 恒等式确定的值 $Z^{{\rm WI}}_A= 1.06(2)$ \cite{wi,gribov}，
见图 \ref{fig:za}。如文献 \cite{wi} 所观察到的，对于 $\beta=6.0$--$6.2$，微扰论给出的 $Z_{A^L}$ 小于一，而 Ward 恒等式
方法给出的值大于一；我们的非微扰结果也提示其值大于一。

\noindent {\bf 赝标密度：}

赝标密度具有标量密度的主要特征，但有两点重要区别：其一它与准 Goldstone 玻色子耦合；其二采用 Clover 作用量时一圈微扰修正
相当大。因此毫不奇怪，如图 \ref{fig:zp} 所示，相应重正化常数的非微扰值远低于微扰论结果。倘若在图 \ref{fig:zp} 中用标准
微扰论而非 boosted 微扰论，偏差还会更大。这一偏差本可以从文献 \cite{wi} 的结果预见到：文中表明，用 Ward 恒等式方法测定的
比值 $Z_S/Z_P$ 显著大于 boosted 微扰论的结果。把这一信息与 $Z_S$ 的非微扰测定和 BPT 结果的一致性相结合，可以断定差别来自
$Z_P$。这也可能是由于 $Z_P$ 有较大的一圈有限修正（$\sim 35 \%$），而且未被单圈 boosted 微扰论计入的高阶项也许很重要。

我们还在图 \ref{fig:zsszp} 中给出作为 $\mu^2$ 函数的 $Z_S/Z_P$。比值 $Z_S/Z_P$ 的行为类似于 $Z_{A^L}$。某些离散化效应似乎比
分别观察 $Z_S$ 与 $Z_P$ 时要小。数值结果定性地支持理论解释：在中间 $\mu$ 范围内该比值几乎保持恒定；在小 $\mu$ 或大 $\mu$ 处，
由于分别为非微扰效应或离散化效应，结果不稳定。在 $\mu^2 a^2 \sim 1 $ 处，非微扰方法得到的该比值与 Ward 恒等式方法测定的
$Z_S/Z_P=1.64(5)$ \cite{wi} 合理一致\footnote{文献 \cite{wi} 用现有统计量一半所发表的值为 $Z_P/Z_S=0.64(2)$，对应于
$Z_S/Z_P=1.56(5)$。}。
